%%% Fiktivní kapitola s ukázkami sazby

\chapter{Teoretická část}
Background, permutace, generátory, satsuma, breakid, detekují symetrie ale není jasné proč fungují dobře generátory

Definice: budeme používat 2D binární matice, seřazení do vektoru a proměnné v řádcích

Definice patternu, odkaz na paper, popis problémů




\section{Základní definice}

Jelikož se budeme zabývat symetriemi, hlavní objekty, na které se budeme soustředit, budou samozřejmě podgrupy symetrické grupy. Prvky, které budeme permutovat, budou proměnné ze zadaných SAT formulí. Ty můžeme reprezentovat jednoduše jako čísla $\{1,...,n\}$, takže symetrickou grupu na těchto prvcích můžeme značit $S_n$.

Budeme předpokládat, že symetrie zadaného problému známe již ze zadání. Je to z toho důvodu, že budeme chtít navázat na jednu konkrétní informaci z nástrojů jako je $BREAKID$ nebo $SATSUMA$ (odkaz na papery). Tyto nástroje dostanou na vstupu SAT formuli a snaží se v ní detekovat určité typy symetrií. Pokud symetrii detekují, tím pádem najdou určitou podgrupu $G \subseteq S_n$, rozbíjí symetrie generováním Lex-Leader klauzulí $x \leq_{\text{lex}} \pi(x)$ pro $\pi$ z množiny generátorů $G$, které získají při detekci symetrie. Zmíněnou informací je, že takto získané klauzule dávají velmi dobré částečné rozbití symetrie, ale ani samotným autorům těchto nástrojů není jasné, proč tomu tak je. V této práci chceme tedy zkoumat, jestli vlastnost, že určitá množina permutací generuje celou podgrupu odpovídající symetrii daného problému , znamená, že s ní dosáhneme lepšího částečného rozbití symetrie. Pro úplnost si na jednoduchém příkladu ukážeme, že množinou generující danou grupu permutací nemusíme dosáhnout úplného rozbití symetrie, což by někoho na první pohled mohlo u generátorů napadnout.

\begin{prikl}
Použijeme příklad již zmíněný v úvodu. Uvažujme tedy problém holubníku s pěti holubníky a jedním holubem. Máme tedy celkem pět proměnných, které budeme nyní pro přehlednost značit $x_1, x_2, x_3, x_4, x_5$. Hodnoty těchto proměnných mohou být pouze $1$ nebo $0$ a hodnota $1$ znamená, že holub je v holubníku odpovídajícímu danému indexu. Z našich pěti proměnných tedy vždy musí právě jedna nabývat hodnoty $1$ a právě čtyři musí nabývat hodnoty $0$. Holubníky jsou si ekvivalentní a tím pádem můžeme proměnné libovolně permutovat. Pracujeme tedy s grupou $S_5$ (kde permutace indexů odpovídá permutaci proměnných s těmito indexy). Pro tu vezmeme množinu generátorů $\{(1,2),(1,5,4,3,2)\}$. Označme si je $\pi_1 := (1,2)$ a $\pi_2 := (1,5,4,3,2)$. Nyní dostaneme dvě podmínky:

\begin{align*}
x = (x_1, x_2, x_3, x_4, x_5) \leq_{\text{lex}} (x_2, x_1, x_3, x_4, x_5) = \pi_1(x) \\
x = (x_1, x_2, x_3, x_4, x_5) \leq_{\text{lex}} (x_2, x_3, x_4, x_5, x_1) = \pi_2(x)
\end{align*}
Dosazením všech možných ohodnocení dostaneme (první sloupec představuje dosazení pro $\pi_1$ a druhý sloupec pro $\pi_2$):
\begin{align*}
(1,0,0,0,0) \leq_{\text{lex}} (0,1,0,0,0) \qquad (1,0,0,0,0) \leq_{\text{lex}} (0,0,0,0,1) \\
(0,1,0,0,0) \leq_{\text{lex}} (1,0,0,0,0) \qquad (0,1,0,0,0) \leq_{\text{lex}} (1,0,0,0,0) \\
(0,0,1,0,0) \leq_{\text{lex}} (0,0,1,0,0) \qquad (0,0,1,0,0) \leq_{\text{lex}} (0,1,0,0,0) \\
(0,0,0,1,0) \leq_{\text{lex}} (0,0,0,1,0) \qquad (0,0,0,1,0) \leq_{\text{lex}} (0,0,1,0,0) \\
(0,0,0,0,1) \leq_{\text{lex}} (0,0,0,0,1) \qquad (0,0,0,0,1) \leq_{\text{lex}} (0,0,0,1,0)
\end{align*}
Pro nás bude vždy platit, že $0 \leq_{\text{lex}} 1$ a že proměnné dříve v pořadí mají větší váhu než proměnné dále. Kanonickým řešením pro nás tedy zde je $(0,0,0,0,1)$. Pokud se nyní podíváme na nerovnosti výše, vídíme, že v obou sloupcích je neplatný první řádek. To znamená, že jsme přidáním podmínek z permutací $\pi_1, \pi_2$ odstranili jedno ekvivalentní řešení z pěti, a to $(1,0,0,0,0)$ (k tomu by stačilo, aby byl neplatný řadek jen v jednom ze sloupců). Pro úplné rozbití symetrie ale potřebujeme odstranit čtyři z pěti těchto ohodnocení.

Analogickým postupem můžeme ověřit, že ale například pro množinu generátorů $\{(1,2),(1,2,3,4,5)\}$ úplné rozbití symetrie dostaneme a zbyde nám jako jediné možné řešení $(0,0,0,0,1)$.
\end{prikl}

V této práci se zaměříme především na problémy, jejichž řešení lze reprezentovat binární maticí. Veškeré postupy je možné přímočaře zobecnit i do vícerozměrných problémů, jen se zpravidla dramaticky zvedne počet proměnných, symetrií a celkově náročnost vyřešení problému. 




\begin{definice}
Nechť $A$ je $m \times n$ matice. Pak $vec(A)$ bude značit vektor délky $m \cdot n$, který získáme konkatenací všech řádků matice za sebe.
\end{definice}

\begin{definice}
Nechť $A, B$ jsou $m \times n$ binární matice a nechť $a := vec(A), b := vec(B)$. Pak budeme uvažovat následující uspořádání $\leq$ definované jako:
\begin{align*}
A \leq B \text{ právě tehdy, když } a \leq_{\text{lex}} b
\end{align*}
\end{definice}

Spíše než abychom ale pracovali s již známými maticemi řešení, budeme chtít lexikograficky menší řešení hledat. Místo binárních matic budou tedy matice obsahovat proměnné a pro každé ohodnocení těchto proměnných dostaneme konkrétní binární matici. Uspořádání $\leq$ z předchozí definice nám poté bude dávat podmínky, které musí ohodnocení proměnných splňovat. 

\begin{prikl}
Vráťme se opět k problému holubníku. Nyní uvažujme, že máme 3 holuby a 3 holubníky. Vytvoříme si matici $H$, ve které budou řádky odpovídat jednotlivým holubům a sloupce holubníkům. Prvky $h_{i,j}$ matice nám pak říkají, jestli $i$-tý holub je v $j$-té holubníku. Základní podmínky, které musíme zakódovat do SAT klauzulí, jsou, že každý holub musí být alespoň v jednom holubníku a v každém holubníku je nejvýše jeden holub. Toto kódování si popíšeme později. Nyní je podstatné, že máme následující matici:

\begin{align*}
H =  \begin{pmatrix}
x_1 & x_2 & x_3 \\
x_4 & x_5 & x_6 \\
x_7 & x_8 & x_9 
\end{pmatrix}
\end{align*}

$H$ musí z výše uvedených podmínek mít v každém řádku alespoň jednu $1$ a v každém sloupci nejvýše jednu $1$. Již víme, že holubníky můžeme libovolně permutovat. Stejně tak zřejmě můžeme permutovat i holuby. V tomto případě tedy pracujeme s grupou $S_3 \times S_3 \subseteq S_9$. Vezmeme-li například pro řádky (holuby) permutaci $\pi_1 = (1,2,3)$ a pro sloupce (holubníky) permutaci $\pi_2 = (2,1,3)$, dostaneme pro $\pi := (\pi_1, \pi_2)$
\begin{align*}
\pi(H) =  \begin{pmatrix}
x_8 & x_9 & x_7 \\
x_2 & x_3 & x_1 \\
x_5 & x_6 & x_4 
\end{pmatrix}
\end{align*}
a přidáme podmínku $H \leq \pi(H)$. Což znamená 
\begin{align*}
vec(H) &\leq_{\text{lex}} vec(\pi(H)) \text{ ,neboli}\\
(x_1, x_2, x_3, x_4, x_5, x_6, x_7, x_8, x_9) &\leq_{\text{lex}} (x_8, x_9, x_7, x_2, x_3, x_1, x_5, x_6, x_4).
\end{align*}
Dále už bychom mohli postupovat analogicky prvnímu příkladu s jedním holubem a pěti holubníky.
\end{prikl}

\section{Kódování podmínek}
V této části se podíváme blíže na to, jak budeme kódovat podmínky $\leq_{\text{lex}}$ jako SAT klauzule. Využijeme k tomu takzvané \emph{vzory} z (odkaz na paper) (anglicky \emph{patterns}). Idea je taková, že pro podmínku $x \leq_{\text{lex}} \pi(x)$ budeme naopak postupně sestavovat podmínku $\pi(x) <_{\text{lex}} x$, kterou poté znegujeme. Postupným sestavováním se myslí, že si v každém kroku zafixujeme pozici ve vektorech, na které se mají poprvé lišit (respektive první pozice, na které je proměnná v $\pi(x)$ ostře menší než proměnná v $x$). To nám zároveň říká, že na všech předchozích pozicích se vektory musí rovnat. Z těchto vlastností pak odvodíme hledaný vzor. 

\begin{definice}
Nechť $a = (a_1,...,a_n)$ a $b = (b_1,...,b_n)$. Pak řekneme, že $a$ je menší než $b$ na $i$-té pozici, pokud
$$(a_1 = b_1) \land ... \land (a_{i-1} = b_{i-1}) \land (a_i < b_i).$$
Tuto vlastnost budeme značit $a <_{\text{lex}}^{i} b$.
\end{definice}

\begin{pozn}
Jelikož pracujeme s binárními proměnnými, $(a_i < b_i)$ znamená jednoduše $(a_i = 0) \land (b_i = 1)$.
\end{pozn}

\begin{pozn}
$$a <_{\text{lex}}^{i} b \Leftrightarrow \bigvee_{i=1}^{n} a <_{\text{lex}}^{i} b$$
\end{pozn}

\begin{definice}
Nechť $\pi$ je permutace a $a = (a_1,...,a_n)$. Pak vzorem $pat_i(\pi)$ rozumíme aplikování rovností a nerovnosti z $\pi(a) <_{\text{lex}}^{i} a$ na vektor $a$.
\end{definice}

\begin{prikl}
Uvažujme vektory z předchozího příkladu a mějme $vec(\pi(H)) <_{\text{lex}}^{i} vec(H)$ pro $i=4$, neboli 
$$(x_8, x_9, x_7, x_2, x_3, x_1, x_5, x_6, x_4) <_{\text{lex}}^{4} (x_1, x_2, x_3, x_4, x_5, x_6, x_7, x_8, x_9).$$
Tím tedy dostáváme podmínky $(x_8 = x_1), (x_9 = x_2), (x_7 = x_3)$ a $(x_2 < x_4)$. To znamená, že $x_2 = 0$, $x_4 = 1$ a zároveň také z druhé rovnosti $x_9 = 0$. Pro odlišení dále označíme $A := x_1$ ($=x_8$) a $B := x_3$ ($=x_7$). Dosazením (aplikováním) těchto hodnot do $vec(H)$ dostaneme 
$$pat_4(\pi) = (A, 0, B, 1, x_5, x_6, B, A, 0).$$
Nyní pokud libovolný vektor $v$ odpovídá tomuto vzoru, tak o něm automaticky víme, že $v <_{\text{lex}} vec(H)$. Jelikož $A, B, x_5, x_6$ mohou být libovolná, pokryli jsme tímto vzorem celkem $2^4 = 16$ různých vektorů.
\end{prikl}






